\section{Funktioner}
\begin{align*}
    f : A \to B \\
    a \mapsto f(a)
\end{align*}
Beskriver funtkionen fra \emph{Definitionsmængden} \(A\) til \emph{Dispositionsmængden} \(B\), hvor et element \(a\) i \(A\) bliver afbildet til \(f(a)\) i \(B\).

\emph{Billedet}/\emph{Værdimængden} af funktionen er mængden af alle \(\{ f(a) | a \in A \}\).

\emph{Billedet} $\subsetneq$ \emph{Dispositionsmængden}.

\subsubsection{Sammensatte funktioner}
Givet to funktioner \(f : A \to B\) og \(g : B \to C\), kan man sammensætte dem til en ny funktion \(g \circ f : A \to C\), hvor \((g \circ f)(a) = g(f(a))\).

\subsubsection{Identitetsfunktionen}
Identitetsfunktionen \(id_A : A \to A\) er defineret som \(id_A(a) = a\) for alle \(a \in A\).

\[f \circ (g \circ h) = (f \circ g) \circ h \kildetag{Lemma 2.2.1}\] 

\subsection{Funktions-egenskaber}
\subsubsection{Injektiv}
Der findes ikke to elementer i definitionsmængden, som bliver afbildet til det samme element i dispositionsmængden.

For all injektive funktioner gælder:
\[
    a_1 < a_2 \Rightarrow (f(a_1) < f(a_2)) \lor (f(a_1) > f(a_2)) \kildetag{Lemma 2.3.1}
\]


\subsubsection{Surjektiv}
Hvert element i dispositionsmængden har mindst ét element i definitionsmængden, som bliver afbildet til det.

\emph{Billedet} = \emph{Dispositionsmængden}.

\subsubsection{Bijektiv}

Funktionen er både injektiv og surjektiv.

Hvis og kun hvis en funktion er bijekitv har den en invers funktion \(f^{-1} : B \to A\), hvor \(f^{-1}(f(a)) = a\) for alle \(a \in A\) og \(f(f^{-1}(b)) = b\) for alle \(b \in B\). \kilde{Lemma 2.2.2}
